\scanpage{22}
\uline{Théorème 2} Le foncteur $\unsure{K-\Pi_K/\Pi_{\mathbb Q}|K\to\Pi_K/\Pi_{\mathbb Q}}$ des catégories des t.f. de $\mathbb Q$ vers les groupoïdes profinis sur $\Pi_{\mathbb Q}$ est \uline{fidèle}.

\margin{En fait, ce théorème n'est pas spécial à $\mathbb Q$ --- il marcherait sur un corps de base quelconque de car $0$. \unsure{La|car la} démon. montre bien que c'est en fait un énoncé de nature géométrique non arithmétique (j'énonce pour \unsure{rendre} la démon. facile \unsure{et|une} arithmétique pure)}

Si deux hom. $K\overset{f,g}{\rightrightarrows}K'$ définissent des hom de groupoïdes \inserted{sur $\Pi_{\mathbb Q}$} isomorphes
\diagram{\begin{tikzcd}[ampersand replacement=\&]
\Pi_{K'} \arrow[rr,shift left,"f^*"] \arrow[rr,shift right,"g^*"'] \arrow[dr,no head,"p'"'] \& \& \Pi_K \arrow[dl,"p"]\\
\& \Pi_{\mathbb Q} \&
\end{tikzcd}}
(i.e. $\exists$ isomorphisme de foncteurs $\alpha:f^*\simeq g^*$, tel que $\forall\bar\eta'\in\operatorname{Ob}\Pi_{K'}$, l'e\deleted{isom}\inserted{carré}
\deleted{correspondant}
\diagram{\begin{tikzcd}[ampersand replacement=\&]
pf^*(\bar\eta') \arrow[r,phantom,"{\overset{p(\alpha)}{\simeq}}" description] \arrow[d,"\wr"'] \& pg^*(\eta) \arrow[d,no head,"\wr"]\\
p'(\bar\eta') \arrow[r,phantom,"\simeq" description] \& p'(\bar\eta')
\end{tikzcd}}
est commutatif) alors $f=g$.

L'hypothèse sur $f,g$ signifie aussi, en termes d'une clôture alg. choisie $\bar K'$ \deleted{\unsure{qui ne|qu'on} \unsure{peut}}\inserted{donnant via $f,g$ deux clôtures alg. de $\unsure{\bar K_\varepsilon|\bar K}$, $\unsure{K|\widetilde K}$ de $K$} de $K'$, que \deleted{les deux hom d'exten}\inserted{l'on peut trouver un isomorphisme entre celles-ci (permettant donc d'identifier $\mathcal E_{\bar K/K}$ et $\mathcal E_{\widetilde K/K}$), de telle façon que les deux hom}
\deleted{\unsure{disons}}
\margin{induisant ``l'identité'' sur les clôtures alg. \deleted{\unsure{et|de}} $\overline{\mathbb Q}$ de $\mathbb Q$ dans $\bar K$ et $\widetilde K$, identifiées à celle dans $\bar K'$}
\[
f^*,g^*:\mathcal E_{\bar K'/K'}\unsure{\rightrightarrows}\mathcal E_{\bar K/K}
\]
sont égaux, \deleted{Il \unsure{s'agit} en \unsure{conclusion}}
% The following cancelled replacement is written above the preceding cancelled run.
\inserted{\deleted{\unsure{C'est} \unsure{sans} doute plus}}
\unsure{Ou bien} en termes d'une clôture alg. $\overline{\mathbb Q}$ fixée, de $\mathbb Q$, en disant que les deux hom $f^*,g^*:\pi_{K'}\rightrightarrows\pi_K$ de groupes profinis extérieurs (avec opérateurs $\Gamma_{\overline{\mathbb Q}/\mathbb Q}$) sont égaux.
\margin{On doit se ramener au cas $k=k'$ --- c'est immédiat en \deleted{\unsure{faisant le chgt de base}}\inserted{\unsure{prenant}} extension finie de $\mathbb Q$}
\margin{21}
